\subsection{按提交顺序重放一次人工准备}

现在把准备过程按地址和顺序展开。假设调试夹具每次最多写四个字节，并能在每笔写后读回。
操作者为本次运行选择编号 17。


处理器保持复位。夹具向 \code{0x001ff030} 写入四字节零，再读回确认。即使 RAM 中其余
区域仍是上一次运行的完整内容，启动程序此刻也不得采用它。

\begin{lstlisting}
write32 0x001ff030 0x00000000
read32  0x001ff030 -> 0x00000000
\end{lstlisting}

若这一步读回失败，不能继续写代码后再“看看是否能运行”。失败已经说明夹具、互连旁路或
RAM 写入路径不可靠，应在机器仍保持复位时诊断。


\begin{longtable}{|L{0.20\linewidth}|L{0.25\linewidth}|L{0.43\linewidth}|}
\hline
\rowcolor{BookTableHead}
地址 & 写入的四个字节 & 读回后解释 \\ \hline
\code{0x00100000} & \code{10 04 00 00} & 清零 \code{r4} \\ \hline
\code{0x00100004} & \code{20 05 00 00} & 从 \code{r0} 指向处装入 \code{r5} \\ \hline
\code{0x00100008} & \code{11 04 05 00} & \code{r4 <- r4+r5} \\ \hline
\code{0x0010000c} & \code{12 00 04 00} & \code{r0 <- r0+4} \\ \hline
\code{0x00100010} & \code{12 01 ff ff} & \code{r1 <- r1-1} \\ \hline
\code{0x00100014} & \code{31 01 fb ff} & 非零时回退五条指令 \\ \hline
\code{0x00100018} & \code{21 04 02 00} & 把 \code{r4} 写到 \code{r2} 指向处 \\ \hline
\code{0x0010001c} & \code{41 00 00 00} & 从任务栈返回 \\ \hline
\end{longtable}

夹具按地址递增写入并读回。读回比较以字节为单位，而不是把四字节先解释成宿主机器的整数，
从而避免夹具所在计算机与教学处理器字节序不同造成误判。


\begin{longtable}{|L{0.22\linewidth}|L{0.24\linewidth}|L{0.42\linewidth}|}
\hline
\rowcolor{BookTableHead}
地址 & 32 位写入值 & 用途 \\ \hline
\code{0x00102000} & \code{0x00000007} & 第一个输入 \(7\) \\ \hline
\code{0x00102004} & \code{0xfffffffe} & 第二个输入 \(-2\) \\ \hline
\code{0x00102008} & \code{0x0000000b} & 第三个输入 \(11\) \\ \hline
\code{0x0010200c} & \code{0x00000005} & 第四个输入 \(5\) \\ \hline
\code{0x00102020} & \code{0xdeadbeef} & 尚未产生结果的哨兵 \\ \hline
\code{0x001efffc} & \code{0x00000300} & 最外层 \code{RET} 的返回地址 \\ \hline
\end{longtable}

结果槽与输入末端之间保留十六字节空隙不是硬件要求，而是为了让本例中的错一位地址更容易
观察。栈下方大部分空间暂未写入，因为任务没有内部调用；不能因此假定那些字节为零。


\begin{lstlisting}
0x001ff000  magic         = bytes "RUN1"
0x001ff004  generation    = 17
0x001ff008  code_base     = 0x00100000
0x001ff00c  code_bytes    = 32
0x001ff010  data_base     = 0x00102000
0x001ff014  data_bytes    = 16
0x001ff018  result_base   = 0x00102020
0x001ff01c  stack_low     = 0x001e0000
0x001ff020  stack_top     = 0x001f0000
0x001ff024  entry         = 0x00100000
0x001ff028  code_sum      = computed checksum
0x001ff02c  data_sum      = computed checksum
0x001ff030  prepared      = 0
\end{lstlisting}

操作者重新从记录读取代码首地址与长度，再按记录指定范围读回代码并复算校验和，而不是只
拿手边原始清单计算。这样可以同时发现“字节写对但记录地址写错”的情况。数据区同理。


全部读回一致后，夹具执行唯一的提交写：

\begin{lstlisting}
write32 0x001ff030 0x00000001
read32  0x001ff030 -> 0x00000001
\end{lstlisting}

随后操作者停止使用夹具写 RAM，再释放复位。若在提交后又修补一个代码字节，原校验和与
“提交后内容不变”的前提同时失效，必须重新清零标志并完整核对。


启动程序的检查顺序可以用一张状态表精确表示。每个失败状态都保留准备编号、阶段和一个
辅助值。

\begin{longtable}{|L{0.17\linewidth}|L{0.30\linewidth}|L{0.41\linewidth}|}
\hline
\rowcolor{BookTableHead}
阶段 & 成功条件 & 失败时 \code{detail} 保存什么 \\ \hline
\code{READY} & \code{prepared=1} & 实际读到的标志 \\ \hline
\code{FORMAT} & \code{magic=RUN1} & 第一个不匹配的字节偏移 \\ \hline
\code{RANGE} & 所有区域范围合法且互不冲突 & 失败字段偏移或区域对 \\ \hline
\code{CODE} & 代码校验和相等 & 重新计算所得值 \\ \hline
\code{DATA} & 输入校验和相等 & 重新计算所得值 \\ \hline
\code{STACK} & 栈顶字为合法返回地址 & 实际返回字 \\ \hline
\code{ENTER} & 再次读取提交标志仍为一 & 第二次读到的值 \\ \hline
\end{longtable}


\begin{longtable}{|L{0.12\linewidth}|L{0.22\linewidth}|L{0.25\linewidth}|L{0.29\linewidth}|}
\hline
\rowcolor{BookTableHead}
快照 & 当前 PC 范围 & 已证明的内容 & 尚未证明的内容 \\ \hline
\(S_0\) & 启动入口 & 复位起点确定 & RAM 是否属于本批次 \\ \hline
\(S_1\) & 记录检查 & 标志、格式、编号正确 & 区域内容完整性 \\ \hline
\(S_2\) & 范围检查 & 地址和长度关系合法 & 实际字节是否匹配 \\ \hline
\(S_3\) & 校验循环 & 代码、输入校验和匹配 & 任务算法是否正确 \\ \hline
\(S_4\) & 交接准备 & 参数、栈和入口候选完整 & 任务是否最终返回 \\ \hline
\(S_5\) & 任务入口 & 入口状态已经提交 & 结果尚未产生 \\ \hline
\end{longtable}

这组快照避免把不同层次的证明混合。校验和匹配只说明扫描到的字节与记录一致，不说明机器
码一定实现求和；算法正确性要由译码表、指令 trace 和循环不变量说明。


假设地址 \code{0x0010000a} 的字节应为 \code{05}，实际读回为 \code{04}。准备记录中的
\code{code\_sum} 仍按正确代码计算。启动程序扫描 32 字节后得到不同总和，于是写：

\begin{lstlisting}
boot_status.phase  = CODE
boot_status.detail = recomputed_sum
boot_status.run_id = 17
processor PC       = checksum_failure_stop
\end{lstlisting}

它不把 PC 改为 \code{0x00100000}，也不保留一部分任务参数作为有效入口状态。操作者保持
或重新拉起复位，清零提交标志，修复字节，再重新扫描整个代码区。只重算出错位置之后的
部分不足以证明此前字节未被同时改变。


本例只有 32 个代码字节、16 个输入字节和一份短记录，操作者尚可逐项核对。若程序扩大到
64 KiB，以每笔四字节写入计算，需要 16384 笔写和同量级读回；任意一次地址抄错都可能
迫使整段重新核对。更换输入还要重复编辑记录和准备编号。

问题不在于操作者“不够仔细”，而在于机器没有承担可机械重复的搬运工作。已经存在的启动
程序能读取 RAM、计算校验和、建立入口状态，却没有获得外部字节。下一章增加的第一个部件
因此应当只解决“怎样让启动程序可靠地接收一串字节”，而不同时引入第二项任务或定时器。
